% !TeX root = ../sections/03_solutions.tex

\begin{solution}
	整数係数行列の単因子を求める。
	ここで，整数行列 \(A\) に対して，ユニモジュラー行列 \(P,Q\) により
	\[
	PAQ=\operatorname{diag}(a_1,\ldots,a_r)
	\]
	とできるとき，\(a_1,\ldots,a_r\) を \(A\) の単因子という。
	ただし
	\[
	a_1\mid a_2\mid \cdots \mid a_r
	\]
	を満たすように取る。
	
	また，\(P,Q\) がユニモジュラーであるとは，
	\[
	\det P=\pm 1,
	\qquad
	\det Q=\pm 1
	\]
	であることをいう。
	
	以下では，\(k\) 次小行列式全体の最大公約数を \(d_k\) と書く。
	このとき，階数が \(r\) の行列について
	\[
	d_1=a_1,
	\qquad
	d_2=a_1a_2,
	\qquad
	d_3=a_1a_2a_3
	\]
	が成り立つ。
	
	\begin{enumerate}
		\item
		\[
		A=
		\begin{pmatrix}
			2&0&0\\
			0&3&0\\
			0&0&4
		\end{pmatrix}
		\]
		とする。
		
		まず，\(1\) 次小行列式は行列の各成分である。
		非零成分は
		\[
		2,\quad 3,\quad 4
		\]
		であるから，
		\[
		d_1=\gcd(2,3,4)=1.
		\]
		したがって，
		\[
		a_1=1
		\]
		である。
		
		次に，\(2\) 次小行列式を考える。
		この行列は対角行列なので，非零となる主な \(2\) 次小行列式は
		\[
		\det
		\begin{pmatrix}
			2&0\\
			0&3
		\end{pmatrix}
		=6,
		\qquad
		\det
		\begin{pmatrix}
			2&0\\
			0&4
		\end{pmatrix}
		=8,
		\qquad
		\det
		\begin{pmatrix}
			3&0\\
			0&4
		\end{pmatrix}
		=12
		\]
		である。その他の \(2\) 次小行列式は \(0\) である。
		したがって，
		\[
		d_2=\gcd(6,8,12)=2.
		\]
		よって，
		\[
		a_1a_2=d_2=2
		\]
		であり，\(a_1=1\) だから，
		\[
		a_2=2
		\]
		である。
		
		最後に，
		\[
		d_3=|\det A|
		\]
		である。
		ここで，
		\[
		\det A=2\cdot 3\cdot 4=24
		\]
		だから，
		\[
		d_3=24.
		\]
		したがって，
		\[
		a_1a_2a_3=d_3=24.
		\]
		すでに
		\[
		a_1=1,
		\qquad
		a_2=2
		\]
		であるから，
		\[
		1\cdot 2\cdot a_3=24
		\]
		となる。よって，
		\[
		a_3=12.
		\]
		
		したがって，単因子は
		\[
		\boxed{1,\ 2,\ 12}
		\]
		である。
		
		実際，次の行列を取る。
		\[
		P=
		\begin{pmatrix}
			1&1&0\\
			-3&-2&1\\
			6&4&-3
		\end{pmatrix},
		\qquad
		Q=
		\begin{pmatrix}
			-1&-3&-6\\
			1&2&4\\
			0&-1&-3
		\end{pmatrix}.
		\]
		このとき，
		\[
		\det P=-1,
		\qquad
		\det Q=-1
		\]
		なので，\(P,Q\) はユニモジュラー行列である。
		
		さらに直接計算すると，
		\[
		PAQ
		=
		\begin{pmatrix}
			1&0&0\\
			0&2&0\\
			0&0&12
		\end{pmatrix}.
		\]
		したがって，
		\[
		\boxed{
			PAQ=\operatorname{diag}(1,2,12)
		}
		\]
		である。
		
		\item
		\[
		A=
		\begin{pmatrix}
			2&1\\
			-4&2
		\end{pmatrix}
		\]
		とする。
		
		まず，\(1\) 次小行列式は各成分である。
		非零成分は
		\[
		2,\quad 1,\quad -4,\quad 2
		\]
		であるから，
		\[
		d_1=\gcd(2,1,-4,2)=1.
		\]
		したがって，
		\[
		a_1=1
		\]
		である。
		
		次に，この行列は \(2\times 2\) 行列なので，
		\[
		d_2=|\det A|
		\]
		である。
		計算すると，
		\[
		\det A
		=
		\det
		\begin{pmatrix}
			2&1\\
			-4&2
		\end{pmatrix}
		=
		2\cdot 2-1\cdot(-4)
		=
		4+4
		=
		8.
		\]
		したがって，
		\[
		d_2=8.
		\]
		ゆえに，
		\[
		a_1a_2=d_2=8.
		\]
		すでに \(a_1=1\) なので，
		\[
		a_2=8.
		\]
		
		したがって，単因子は
		\[
		\boxed{1,\ 8}
		\]
		である。
		
		実際，次の行列を取る。
		\[
		P=
		\begin{pmatrix}
			1&0\\
			2&-1
		\end{pmatrix},
		\qquad
		Q=
		\begin{pmatrix}
			0&1\\
			1&-2
		\end{pmatrix}.
		\]
		このとき，
		\[
		\det P=-1,
		\qquad
		\det Q=-1
		\]
		であるから，\(P,Q\) はユニモジュラー行列である。
		
		さらに，
		\[
		PAQ
		=
		\begin{pmatrix}
			1&0\\
			0&8
		\end{pmatrix}
		\]
		である。
		したがって，
		\[
		\boxed{
			PAQ=\operatorname{diag}(1,8)
		}
		\]
		である。
		
		\item
		\[
		A=
		\begin{pmatrix}
			2&1&-1\\
			-2&-3&1\\
			-2&-1&3
		\end{pmatrix}
		\]
		とする。
		
		まず，\(1\) 次小行列式は各成分である。
		この行列には成分 \(1\) が含まれているので，
		\[
		d_1=1.
		\]
		したがって，
		\[
		a_1=1
		\]
		である。
		
		次に，\(2\) 次小行列式を調べる。
		たとえば，第 \(1,2\) 行，第 \(1,2\) 列からなる小行列式は
		\[
		\det
		\begin{pmatrix}
			2&1\\
			-2&-3
		\end{pmatrix}
		=
		2\cdot(-3)-1\cdot(-2)
		=
		-6+2
		=
		-4
		\]
		である。
		また，第 \(1,2\) 行，第 \(1,3\) 列からなる小行列式は
		\[
		\det
		\begin{pmatrix}
			2&-1\\
			-2&1
		\end{pmatrix}
		=
		2\cdot 1-(-1)(-2)
		=
		2-2
		=
		0
		\]
		である。
		さらに，第 \(1,2\) 行，第 \(2,3\) 列からなる小行列式は
		\[
		\det
		\begin{pmatrix}
			1&-1\\
			-3&1
		\end{pmatrix}
		=
		1\cdot 1-(-1)(-3)
		=
		1-3
		=
		-2
		\]
		である。
		
		すべての \(2\) 次小行列式の最大公約数は
		\[
		d_2=2
		\]
		である。
		したがって，
		\[
		a_1a_2=d_2=2.
		\]
		\(a_1=1\) なので，
		\[
		a_2=2
		\]
		である。
		
		最後に，
		\[
		d_3=|\det A|
		\]
		を求める。
		\[
		\det A
		=
		\det
		\begin{pmatrix}
			2&1&-1\\
			-2&-3&1\\
			-2&-1&3
		\end{pmatrix}.
		\]
		第 \(1\) 行で展開すると，
		\[
		\begin{aligned}
			\det A
			&=
			2
			\det
			\begin{pmatrix}
				-3&1\\
				-1&3
			\end{pmatrix}
			-
			1
			\det
			\begin{pmatrix}
				-2&1\\
				-2&3
			\end{pmatrix}
			+
			(-1)
			\det
			\begin{pmatrix}
				-2&-3\\
				-2&-1
			\end{pmatrix}
			\\
			&=
			2\{(-3)\cdot 3-1\cdot(-1)\}
			-
			\{(-2)\cdot 3-1\cdot(-2)\}
			-
			\{(-2)(-1)-(-3)(-2)\}
			\\
			&=
			2(-9+1)
			-
			(-6+2)
			-
			(2-6)
			\\
			&=
			2(-8)-(-4)-(-4)
			\\
			&=
			-16+4+4
			\\
			&=
			-8.
		\end{aligned}
		\]
		よって，
		\[
		d_3=|\det A|=8.
		\]
		したがって，
		\[
		a_1a_2a_3=d_3=8.
		\]
		すでに
		\[
		a_1=1,
		\qquad
		a_2=2
		\]
		であるから，
		\[
		1\cdot 2\cdot a_3=8
		\]
		となる。ゆえに，
		\[
		a_3=4.
		\]
		
		したがって，単因子は
		\[
		\boxed{1,\ 2,\ 4}
		\]
		である。
		
		実際，次の行列を取る。
		\[
		P=
		\begin{pmatrix}
			1&0&0\\
			-3&-1&0\\
			4&1&1
		\end{pmatrix},
		\qquad
		Q=
		\begin{pmatrix}
			0&0&1\\
			1&1&0\\
			0&1&2
		\end{pmatrix}.
		\]
		このとき，
		\[
		\det P=-1,
		\qquad
		\det Q=1
		\]
		であるから，\(P,Q\) はユニモジュラー行列である。
		
		さらに直接計算すると，
		\[
		PAQ
		=
		\begin{pmatrix}
			1&0&0\\
			0&2&0\\
			0&0&4
		\end{pmatrix}.
		\]
		したがって，
		\[
		\boxed{
			PAQ=\operatorname{diag}(1,2,4)
		}
		\]
		である。
	\end{enumerate}
\end{solution}